\exercisesection{5}

\begin{enumerate}
\def\labelenumi{\arabic{enumi}.}
\tightlist
\item
令 A 为交换环，\(\mathfrak{I}\) 为 A 的一个理想。称 A-模 M 关于 \(\mathfrak{I}\) 是绝对 Hausdorff 的，是指对于每个有限生成的 A-模 N，A-模 N \(\otimes_{\mathrm{A}}\) M 关于 \(\mathfrak{I}\) -进拓扑都是 Hausdorff 的；绝对 Hausdorff 的 A-模称为理想 Hausdorff 模。
\end{enumerate}

\begin{enumerate}
\def\labelenumi{(\alph{enumi})}
\item
M 关于 \(\mathfrak{I}\) 是绝对 Hausdorff 的充要条件是：对于每个有限生成的 A-模 N 及 N 的每个子模 \(\mathbf{N}'\)，\(\operatorname{Im}(\mathbf{N}' \otimes_{\mathbf{A}} \mathbf{M})\) 在 \(\mathbf{N} \otimes_{\mathbf{A}} \mathbf{M}\) 的 \(\mathfrak{I}\) -进拓扑下是闭的。
\item
  设 B 是交换的 A-代数，\(\mathfrak{L}\) 是 B 的一个包含 \(\mathfrak{I}B\) 的理想。如果 M 是关于 \(\mathfrak{L}\) 绝对 Hausdorff 的 B-模，则 M 是关于 \(\mathfrak{I}\) 绝对 Hausdorff 的 A-模。
\end{enumerate}

\begin{enumerate}
\def\labelenumi{\arabic{enumi}.}
\setcounter{enumi}{1}
\item
  证明每个关于 p-进拓扑为 Hausdorff 的 Z-模（p 为素数）都是理想 Hausdorff 模；但给出一个关于 p 为 Hausdorff 而非绝对 Hausdorff 的有限生成 Z-模的例子（使用习题 1 (a)）。
\item
用第 3 节习题 11 的记号。设 N 是 \(\hat{\mathbf{E}}\) 的子模，它是 \(\hat{\mathbf{E}}\) 中由 \(\hat{\mathbf{E}}\) 的子模所生成的向量 \(pe_{2n-1} - p^n e_{2n} (n \geqslant 1)\) 的闭包；令 M = \(\hat{\mathbf{E}}/\mathbf{N}\)。证明在 \(p\) -进拓扑下，M 的子模 \(p\mathbf{M}\) 在 M 中不是闭的，并由此推出 M 是理想 Hausdorff 的 \(\mathbf{Z}_p\)-模，但关于 \(p\) 不是绝对 Hausdorff 的。
\end{enumerate}

¶ 4. 设 A 为交换环，\(\mathfrak{I}\) 为 A 的一个理想，且 S = 1 + \(\mathfrak{I}\)。证明，如果 A-模 M 关于 \(\mathfrak{I}\) 是绝对 Hausdorff 的，则 S \(^{-1}\) M = M，换言之，对于所有 \(a \in\) S，\(x \mapsto ax\) 是 M 到自身的双射。（先证明 M 关于 \(\mathfrak{I}\)-进拓扑为 Hausdorff 的假设蕴含 \(x \mapsto ax\) 是单射；然后证明 M 的子模 aM 关于 \(\mathfrak{I}\)-进拓扑在 M 中稠密，并使用习题 1 (a)。）

\begin{enumerate}
\def\labelenumi{\arabic{enumi}.}
\setcounter{enumi}{4}
\tightlist
\item
  取 \(\mathbf{A} = \mathbf{Z}\) 及 \(\mathfrak{I} = p\mathbf{Z}\)，其中 \(p\) 为素数。
\end{enumerate}

\begin{enumerate}
\def\labelenumi{(\alph{enumi})}
\item
  证明，如果 \(q\) 是不同于 \(p\) 的素数，则 \(\mathbf{Z}\)-模 \(\mathbf{Z} / q\mathbf{Z}\) 满足定理 1 的性质 (ii)，但不满足性质 (i)。
\item
  证明 \(\mathbf{Z}\)-模 \(\mathbf{Q} / \mathbf{Z}\) 满足定理 1 的性质 (iv)，但不满足性质 (ii)。
\end{enumerate}

\begin{enumerate}
\def\labelenumi{\arabic{enumi}.}
\setcounter{enumi}{5}
\tightlist
\item
  设 A 是交换诺特环，\(\mathfrak{I}\) 是 A 的一个理想，M 是一个 A-模。证明定理 1 的条件 (v) 等价于下述条件：
\end{enumerate}

(v') 对于每个有限生成的 A-模 N 及 N 的每个子模 N'，典范映射 \((\mathbf{M} \otimes_{\mathbf{A}} \mathbf{N}')^{\wedge} \to (\mathbf{M} \otimes_{\mathbf{A}} \mathbf{N})^{\wedge}\) 是单射（其中两边分别是 M \(\mathfrak{I}\) 和 M \(\otimes_{A} N'\) 上的 \(\otimes_{A} N\) -进拓扑所对应的 Hausdorff 完备化）。（证明 (v) 蕴含 (v') 时，仿照定理 1 的证明中部分 (E) 的论证。为看出 (v') 蕴含 (v)，考虑被 \(\mathfrak{I}\) 的某个幂湮灭的 A-模 N。）

¶ 7. 设\((A_{\lambda}, f_{\mu\lambda})\)是诺特局部环的有向直接系统；若\(m_{\lambda}\)是\(A_{\lambda}\)的极大理想，假设对于\(\lambda \leqslant \mu\)，有\(m_{\mu} = m_{\lambda}A_{\mu}\)且\(A_{\mu}\)是平坦的\(A_{\lambda}\)。证明\(A = \lim A_{\lambda}\)是诺特环，并且对于所有\(A_{\lambda}\)都是平坦的\(\lambda\)。（若\(m = m_{\lambda}A\)是\(A\)的极大理想（第二章，第 3 节，习题 16），证明在\(A\)上 m-进拓扑是 Hausdorff 的，注意到\(A\)是忠实平坦的\(A_{\lambda}\)对于所有\(\lambda\)。（然后证明，若\(\hat{A}\)是\(A\)关于 m-进拓扑的完备化，则\(\hat{A}\)是诺特环；利用 § 2，第 10 节定理 2 的推论 5；最后证明，对于所有\(\lambda\)，\(\hat{A}\)都是平坦的\(A_{\lambda}\)-模，利用第 2 节定理 1 和第 3 节命题 2。）

¶ 8. 设 A 是诺特局部环，m 是其极大理想，k = A/m 是其剩余域。设 K 是 k 的一个扩域；证明存在从 A 到某个诺特局部环 B 的局部同态，使得 B/mB 同构于 K，且 B 是平坦的 A-模。（先考虑 \(\mathbf{K} = k(t)\) 的情形，根据 t 在 k 上是代数的还是超越的分两种情况；再考虑 K 的包含 k 的子域族 \((\mathrm{K}_{\lambda})\)，该族按包含关系良序，且若 \(K_{\lambda}\) 有前驱 \(\mathrm{K}_{\mu}, \mathrm{K}_{\lambda} = \mathrm{K}_{\mu}(t_{\mu})\)，其中 \(t_{\mu} \in K_{\mu}\)；最后应用习题 7。）

第 IV 章(*)

